% ============================================================
% FASES 2, 3 Y 4 DEL HITO 5
% Deteccion de deriva + Alertas + Simulacion + Prometheus/Grafana
% ============================================================

\subsection{Fase 2 -- Deteccion de Deriva de Datos (Data Drift)}

\subsubsection{Justificacion del problema}

En la industria vitivinicola, las caracteristicas fisicoquimicas del vino varian entre cosechas debido a factores climaticos, de suelo y tecnicas de produccion. Un modelo entrenado con datos de una campana puede degradarse cuando se aplica a vinos de una cosecha posterior cuya distribucion de variables ha cambiado. Esta situacion se conoce como \textit{covariate shift} o deriva de datos: la distribucion de las variables de entrada $P(\mathbf{X})$ cambia respecto al periodo de entrenamiento, mientras que la relacion $P(Y|\mathbf{X})$ se mantiene teoricamente constante.

Siguiendo la taxonomia del profesorado, se distinguen tres tipos de deriva:
\begin{itemize}
    \item \textbf{Covariate shift:} cambia la distribucion de las entradas (lo que detectamos).
    \item \textbf{Label shift:} cambia la distribucion de las clases de salida.
    \item \textbf{Concept drift:} cambian las reglas que relacionan entrada y salida.
\end{itemize}

Para VinOps se ha seleccionado la deteccion de \textbf{covariate shift} porque:
\begin{enumerate}
    \item Es la mas frecuente en el dominio vitivinicola (cambios de cosecha).
    \item No requiere ground truth (etiqueta real) para su deteccion, a diferencia del concept drift.
    \item Es la mas sencilla de implementar y explicar, siguiendo la recomendacion del profesorado.
\end{enumerate}

\subsubsection{Dise\~{n}o del detector}

Se ha implementado un detector de drift con las siguientes caracteristicas:

\paragraph{Ventana deslizante.} Dado el volumen reducido del dataset (178 instancias), se configura una ventana de 20 muestras con evaluacion cada 10 predicciones. Esta configuracion equilibra:
\begin{itemize}
    \item \textbf{Sensibilidad:} 20 muestras son suficientes para estimar 13 variables.
    \item \textbf{Latencia:} evaluar cada 10 predicciones evita sobrecarga computacional.
    \item \textbf{Memoria:} la ventana se mantiene en memoria (lista) sin persistencia externa.
\end{itemize}

\paragraph{Tecnicas implementadas.} Se han seleccionado tres tecnicas complementarias, todas basadas en \texttt{scipy} para mantener el contenedor ligero:

\begin{table}[htbp]
\centering
\caption{Tecnicas de deteccion de drift implementadas}
\label{tab:tecnicas-drift}
\begin{tabular}{@{}llp{5cm}l@{}}
\toprule
\textbf{Tecnica} & \textbf{Tipo} & \textbf{Umbrales} & \textbf{Justificacion} \\
\midrule
Kolmogorov-Smirnov (KS) & Univariado & $p < 0.05$ por variable & Prueba no parametrica estandar para comparar distribuciones continuas \\
Population Stability Index (PSI) & Univariado & $PSI > 0.25$ (alerta), $> 0.1$ (monitor) & Metrica estandar en scoring y riesgo crediticio; facil interpretacion por umbrales \\
Wasserstein (media estandarizada) & Multivariado & $W > 2.0$ (alerta), $> 1.0$ (monitor) & Detecta cambios conjuntos que las tecnicas univariadas podrian omitir \\
\bottomrule
\end{tabular}
\end{table}

\paragraph{KS test.} Contrasta la hipotesis nula de que dos muestras provienen de la misma distribucion continua. Se aplica por variable entre el dataset de referencia (entrenamiento) y la ventana deslizante de produccion. Un $p$-value $< 0.05$ indica que la distribucion de esa variable ha cambiado significativamente.

\paragraph{PSI.} Mide la magnitud del desplazamiento de la distribucion dividiendo en 10 bins y comparando proporciones. Los umbrales son:
\begin{itemize}
    \item $PSI < 0.1$: sin cambio significativo.
    \item $0.1 \leq PSI < 0.25$: monitorizar.
    \item $PSI \geq 0.25$: drift detectado (alerta critica).
\end{itemize}

\paragraph{Wasserstein multivariado.} Calcula la distancia de Wasserstein por variable, normalizada por la desviacion estandar de referencia, y promedia. Detecta cambios conjuntos que las tecnicas univariadas podrian omitir por efecto de compensacion.

\subsubsection{Decision global de drift}

El sistema combina las tres tecnicas mediante una regla de decision jerarquica:

\begin{equation}
    \text{status} = \begin{cases}
        \text{DRIFT\_DETECTED} & \text{si } N_{KS} \geq 3 \lor N_{PSI} \geq 2 \lor W > 2.0 \\
        \text{MONITOR} & \text{si } N_{KS} \geq 1 \lor N_{PSI} \geq 1 \lor W > 1.0 \\
        \text{OK} & \text{en otro caso}
    \end{cases}
\end{equation}

Donde $N_{KS}$ es el numero de variables con $p < 0.05$ y $N_{PSI}$ el numero de variables con $PSI > 0.25$.

\subsubsection{Endpoint /drift}

Devuelve un JSON con la siguiente estructura:
\begin{itemize}
    \item \texttt{status}: OK / MONITOR / DRIFT\_DETECTED.
    \item \texttt{recommendation}: OK / MONITOR / RETRAIN.
    \item \texttt{ks\_test}: p-values por variable, variables alertadas, numero de alertas.
    \item \texttt{psi}: scores por variable, variables alertadas, numero de alertas.
    \item \texttt{wasserstein}: distancia multivariada y umbrales.
\end{itemize}

\subsection{Fase 3 -- Sistema de Alertas}

\subsubsection{Seleccion del canal: Telegram Bot}

Se ha seleccionado Telegram como canal de notificaciones por las siguientes razones tecnicas y operativas:

\begin{enumerate}
    \item \textbf{Simplicidad tecnica:} la API de Telegram es HTTP REST sin autenticacion compleja (unicamente un token). No requiere servidor SMTP, conectores de Teams ni permisos de tenant.
    \item \textbf{Visibilidad:} el grupo de alertas incluye al equipo tecnico y al stakeholder (director tecnico de la bodega), permitiendo que todas las partes interesadas vean las notificaciones.
    \item \textbf{Demo en vivo:} durante la presentacion, el tribunal puede ver las alertas en tiempo real en el chat del grupo.
    \item \textbf{Sin infraestructura propia:} no requiere mantener un servidor de correo ni configurar webhooks externos.
\end{enumerate}

\subsubsection{Throttling}

Para evitar spam en escenarios de degradacion continua, se implementa un mecanismo de \textit{throttling} de 60 segundos por tipo de alerta. Esto garantiza que el canal no se sature si el sistema detecta drift persistente durante varios minutos.

\subsubsection{Tipos de alerta implementados}

\begin{table}[htbp]
\centering
\caption{Tipos de alerta implementados en VinOps}
\label{tab:alertas}
\begin{tabular}{@{}llp{5cm}l@{}}
\toprule
\textbf{Tipo} & \textbf{Trigger} & \textbf{Mensaje} & \textbf{Clase} \\
\midrule
Drift & $N_{KS} \geq 3$ o $N_{PSI} \geq 2$ o $W > 2.0$ & Deriva detectada, reentrenar & Modelo \\
Anomalia & Tasa de anomalias $> 25\%$ en ventana & Tasa de anomalias elevada & Modelo \\
Confianza baja & Confianza media $< 0.65$ en ventana & Confianza media baja & Modelo \\
Operativa & CPU $> 80\%$ o RAM $> 300$ MB & Alerta de recurso & Operativa \\
Simulacion & Endpoint /simulate/* invocado & Simulacion completada & Test \\
\bottomrule
\end{tabular}
\end{table}

Cada alerta incluye detalles tecnicos (valores, umbrales, variables afectadas) y se envia con formato Markdown para facilitar la lectura en dispositivos moviles.

\subsubsection{Simulacion de entorno anomalo}

Se han implementado dos endpoints de simulacion para validar el sistema sin datos reales de produccion:

\paragraph{/simulate/anomaly.} Genera $N$ muestras con valores extremos ($\pm 2.5\sigma$ respecto a la media de referencia) pero dentro de los rangos fisicos del vino. Permite verificar que el detector de anomalias y el sistema de alertas responden correctamente ante muestras atipicas.

\paragraph{/simulate/drift.} Genera $N$ muestras donde una variable especifica esta desplazada $+k\sigma$ respecto a la distribucion de entrenamiento, manteniendo el resto de variables normales. Simula un cambio real en la linea de produccion (ej. uva de otra zona con mayor contenido alcoholico).

Ambas simulaciones inyectan las muestras en la ventana de drift, evaluan el estado y envian una alerta resumen al canal de Telegram.

\subsection{Fase 4 -- Prometheus + Grafana}

\subsubsection{Arquitectura de monitorizacion}

Se ha desplegado una arquitectura de tres capas para la visualizacion de metricas:

\begin{enumerate}
    \item \textbf{Capa de emision:} la API de inferencia expone metricas en formato Prometheus mediante el endpoint \texttt{GET /metrics}.
    \item \textbf{Capa de recoleccion:} Prometheus (puerto 9090) scrapea las metricas cada 5 segundos desde el contenedor de la API.
    \item \textbf{Capa de visualizacion:} Grafana (puerto 3000) consume los datos de Prometheus y los presenta en un dashboard preconfigurado.
\end{enumerate}

\subsubsection{Metricas expuestas}

Las metricas operativas y de modelo se exponen mediante la libreria \texttt{prometheus\_client}:

\begin{table}[htbp]
\centering
\caption{Metricas Prometheus instrumentadas}
\label{tab:metricas-prometheus}
\begin{tabular}{@{}llp{5cm}@{}}
\toprule
\textbf{Metrica} & \textbf{Tipo} & \textbf{Descripcion} \\
\midrule
\texttt{vinops\_requests\_total} & Counter & Peticiones HTTP por metodo, endpoint y status \\
\texttt{vinops\_request\_latency\_seconds} & Histogram & Latencia de peticiones (buckets configurados) \\
\texttt{vinops\_cpu\_usage\_percent} & Gauge & Uso de CPU del proceso \\
\texttt{vinops\_memory\_usage\_bytes} & Gauge & Consumo de RAM residente \\
\texttt{vinops\_prediction\_confidence} & Histogram & Confianza maxima de la prediccion \\
\texttt{vinops\_prediction\_entropy} & Histogram & Entropia de la distribucion de clases \\
\texttt{vinops\_predictions\_by\_class\_total} & Counter & Predicciones acumuladas por cultivar \\
\texttt{vinops\_anomaly\_flags\_total} & Counter & Muestras marcadas como anomalas \\
\texttt{vinops\_low\_confidence\_total} & Counter & Predicciones con confianza $< 0.70$ \\
\bottomrule
\end{tabular}
\end{table}

\subsubsection{Dashboard de Grafana}

El dashboard ``VinOps - Monitorizacion'' incluye 8 paneles preconfigurados:

\begin{table}[htbp]
\centering
\caption{Paneles del dashboard de Grafana}
\label{tab:grafana-panels}
\begin{tabular}{@{}llp{5cm}@{}}
\toprule
\textbf{Panel} & \textbf{Query Prometheus} & \textbf{Que muestra} \\
\midrule
CPU Usage & \texttt{vinops\_cpu\_usage\_percent} & Uso de CPU del proceso en tiempo real \\
Memory Usage & \texttt{vinops\_memory\_usage\_bytes} & RAM consumida \\
Requests per Second & \texttt{rate(vinops\_requests\_total[1m])} & Throughput de la API \\
Request Latency & \texttt{histogram\_quantile(0.95, ...)} & Latencia p50 y p95 en ms \\
Avg Model Confidence & \texttt{avg(vinops\_prediction\_confidence)} & Confianza media del modelo \\
Predictions by Class & \texttt{rate(vinops\_predictions\_by\_class\_total)} & Distribucion de cultivares predichos \\
Anomalies & Low Confidence & \texttt{rate(vinops\_anomaly\_flags\_total)} & Tasa de anomalias detectadas \\
Prediction Entropy & \texttt{avg(vinops\_prediction\_entropy)} & Entropia media (confusion del modelo) \\
\bottomrule
\end{tabular}
\end{table}

\subsubsection{Decisiones de diseno}

\begin{enumerate}
    \item \textbf{No se usa Alibi Detect ni NannyML:} aunque son herramientas de referencia del estado del arte, se opto por implementacion propia con \texttt{scipy} para mantener el contenedor ligero y evitar dependencias pesadas en un entorno academico. En produccion real se migraria a dichas herramientas.
    \item \textbf{Token en .env, no en codigo:} la seguridad del bot se gestiona via variables de entorno, con \texttt{.env} en \texttt{.gitignore}.
    \item \textbf{Metricas operativas y de modelo separadas:} siguiendo la taxonomia del profesor, las alertas operativas (CPU/RAM) y las de modelo (drift/confianza) tienen tipos distintos y umbrales independientes.
    \item \textbf{Alertas desde Python, no desde Grafana:} se opto por envio directo desde la API para tener control programatico del throttling y del contenido, dejando Grafana exclusivamente para visualizacion.
\end{enumerate}

\subsection{Evidencias de funcionamiento}

\subsubsection{Pruebas ejecutadas}

Se han ejecutado las siguientes pruebas sobre el sistema desplegado:

\begin{enumerate}
    \item \textbf{20 predicciones normales:} ventana llenada sin errores. Drift detectado (13/13 variables alertadas por KS y PSI) debido a la repeticion de la misma muestra, lo cual es esperado.
    \item \textbf{Simulacion anomala:} 5 muestras extremas generadas. Alerta enviada a Telegram.
    \item \textbf{Simulacion de drift:} 20 muestras con Alcohol desplazada $+3\sigma$. PSI de Alcohol = 15.4 (umbral 0.25). Alerta enviada a Telegram.
    \item \textbf{Prometheus targets:} estado UP confirmado para el job \texttt{vinops}.
    \item \textbf{Grafana dashboard:} 8 paneles visibles con datos en tiempo real.
\end{enumerate}

\subsubsection{Limitaciones conocidas}

\begin{itemize}
    \item El F1-Score, AUC y precision no se calculan online porque la etiqueta real del cultivar no esta disponible en el momento de la prediccion. Se utilizan metricas proxy (confianza, entropia) y se documenta que el calculo offline requiere el feedback loop (endpoint /feedback).
    \item El uso de CPU aparece como 0\% en Grafana porque el modelo Random Forest es ligero (178 muestras, 13 features) y una prediccion consume microsegundos de CPU. Para observar CPU $> 0\%$ se requeriria saturar la API con cientos de peticiones por segundo.
\end{itemize}
